%! Author = Admin
%! Date = 17.04.2026
\section{Заключение}

В ходе работы была построена связная математическая схема управления квадрокоптером в задаче движения вдоль пространственной траектории.
Она включает нелинейную 16-мерную модель объекта, геометрическое описание кривой, наблюдатель ближайшей точки, расширенный наблюдатель производных и закон управления по выходу с динамическим расширением.

На основе анализа кода и диссертации~\cite{kim_algorithms} можно сделать следующие выводы:
\begin{enumerate}
    \item базовый аналитический регулятор обеспечивает устойчивое слежение по выходу при корректном выборе коэффициентов $\kappa$, $\gamma$ и шага интегрирования;
    \item работоспособность контроллера подтверждается набором автоматизированных тестов на прямых, окружностях и спиралях;
    \item сформированный датасет связывает локальную геометрию траектории и текущее состояние дрона с максимально допустимой устойчивой скоростью движения;
    \item полносвязная нейронная сеть решает задачу адаптивного подбора $V^{\ast}$ и может использоваться как интеллектуальная надстройка над аналитическим регулятором;
    \item визуализация траекторий, ошибок и скоростей позволяет сравнивать качество слежения и выявлять режимы, требующие перенастройки параметров.
\end{enumerate}\hfill\break
Практически работа демонстрирует комбинированный подход, в котором строгая математическая модель и регулятор по выходу обеспечивают устойчивость, а методы машинного обучения улучшают эксплуатационные характеристики системы.
Такой подход хорошо согласуется с задачей интеллектуального управления движением нелинейного объекта в трехмерном пространстве.